\documentclass[10pt]{beamer}

% --- Paquetes ---
\usetheme{metropolis} % Requiere tener instalado mtheme o puedes cambiarlo por \usetheme{Madrid}
\usepackage[utf8]{inputenc}
\usepackage[spanish]{babel}
\usepackage{amsmath}
\usepackage{amsfonts}
\usepackage{amssymb}
\usepackage{graphicx}
\usepackage{booktabs}
\usepackage{hyperref}

% --- Información de Portada ---
\title{Detección de enfermedades en hojas usando visión por computador}
\subtitle{Práctica II: Procesamiento de imágenes y descriptores clásicos}
\author{Grupo N° 4}
\institute{Universidad de Ingeniería y Tecnología (UTEC) \\ Visión por Computador}
\date{22 de abril de 2026}

\begin{document}

% Slide 1: Portada
\maketitle

% Slide 2: Objetivo
\begin{frame}{Objetivo del Proyecto}
    Diseñar e implementar un sistema de análisis de imágenes que permita:
    \begin{itemize}
                \item \textbf{Segmentar} hojas utilizando técnicas de umbralización.
                \item \textbf{Extraer} características visuales mediante descriptores clásicos.
                \item \textbf{Clasificar} enfermedades utilizando el conjunto de datos \textit{PlantVillage}.
    \end{itemize}
\end{frame}

% Slide 3: Pipeline de Procesamiento
\begin{frame}{Pipeline de Procesamiento}
    El flujo de trabajo se divide en siete etapas fundamentales:
    \begin{enumerate}
        \item Preprocesamiento (Escala de grises y filtrado).
        \item Segmentación (Global y Otsu).
        \item Operaciones morfológicas.
        \item Detección de bordes (Canny/Sobel).
        \item Detección de esquinas (Harris/Shi-Tomasi).
        \item Extracción de descriptores (LBP, HOG, BVoW, Custom).
        \item Clasificación (k-NN, SVM, RF).
    \end{enumerate}
\end{frame}

\begin{frame}{Bitácora de Experimentación}
    Debido a la alta cantidad de clases en \textit{PlantVillage}, se realizó un \textbf{Educated Guess} inicial sobre una tarea binaria (Tomato: Healthy vs Unhealthy), donde se generalizó cualquier enfermedad como \textit{non healthy} y luego sobre la tarea original (Tomato: 10 clases originales); para fijar parámetros base:
    \newline
    \begin{columns}
        \begin{column}{0.45\textwidth}
            \textbf{Parámetros Fijos Iniciales:}
            \begin{itemize}
                \item \small Preproc: Gaussian
                \item \small Seg: Otsu
                \item \small Edge: Canny
                \item \small Corner: Shi-Tomasi
                \item \small Feature: HOG
            \end{itemize}
        \end{column}
        \begin{column}{0.55\textwidth}
            \textbf{Variantes de Morfología:}
            \begin{enumerate}
                \item \small Opening only.
                \item \small Erosion $\rightarrow$ Dilation.
                \item \small Opening $\rightarrow$ Closing.
            \end{enumerate}
        \end{column}
    \end{columns}
\end{frame}

\begin{frame}{Resultados de experimentación (Tomato - Binary)}
    Para validar la capacidad discriminativa del pipeline, se realizó una \textbf{binarización estratégica} del dataset de Tomate (9 patologías vs. Healthy) para fijar los mejores parámetros de preprocesamiento.
        
        \begin{table}[]
            \centering
            \begin{tabular}{llccc}
                \toprule
                \textbf{Rank} & \textbf{Configuración Morfología} & \textbf{Modelo} & \textbf{Acc} & \textbf{F1-Score} \\
                \midrule
                \textbf{01} & \textbf{Opening + Closing} & \textbf{SVM} & \textbf{0.8579} & \textbf{0.8578} \\
                02 & Opening only & SVM & 0.8434 & 0.8429 \\
                03 & Erosion + Dilation & SVM & 0.8434 & 0.8429 \\
                04 & Opening only & Rand. F & 0.7975 & 0.7962 \\
                05 & Erosion + Dilation & Rand. F & 0.7975 & 0.7962 \\
                \bottomrule
            \end{tabular}
        \end{table}
        
        \textbf{Análisis de resultados:} 
        \begin{itemize}
            \item El refinamiento mediante \textbf{Opening + Closing} demostró ser superior al eliminar ruido residual y reconstruir la estructura interna de las lesiones generalizadas, con SVM.
        \end{itemize}
\end{frame}


\begin{frame}{Resultados de experimentación (Tomato - Multiclass)}
    Al evaluar las 10 clases originales, se confirma que la configuración ganadora en el escenario binario mantiene su dominancia, a pesar del incremento en la complejidad del problema.
        
        \begin{table}[]
            \centering
            \begin{tabular}{llccc}
                \toprule
                \textbf{Rank} & \textbf{Configuración Morfología} & \textbf{Modelo} & \textbf{Acc} & \textbf{F1-Score} \\
                \midrule
                \textbf{01} & \textbf{Opening + Closing} & \textbf{SVM} & \textbf{0.4194} & \textbf{0.4054} \\
                02 & Opening only & SVM & 0.3941 & 0.3784 \\
                03 & Erosion + Dilation & SVM & 0.3941 & 0.3784 \\
                04 & Opening only & Rand. F & 0.3392 & 0.3146 \\
                05 & Erosion + Dilation & Rand. F & 0.3392 & 0.3146 \\
                06 & Opening + Closing & Rand. F & 0.3419 & 0.3127 \\
                \bottomrule
            \end{tabular}
        \end{table}
        
        \textbf{Análisis de resultados:} 
        \begin{itemize}
            \item La caída del F1-score ($\approx$ 0.40) es esperada debido a la \textbf{similitud morfológica} entre diferentes enfermedades que los descriptores clásicos (HOG/LBP) no logran explicar decentemente. \textit{Aun asi, el F1-score es mayor al baseline de $0.1$}
        \end{itemize}
\end{frame}

\begin{frame}{Conclusiones de experimentación}
   El pipeline es bueno para obtener cosas como \textbf{la presencia de enfermedad o no (binary clasification)}, pero no es suficientemente bueno como para obtener el detalle de que enfermedad posee \textit{(Limite de los descriptores)}.
\end{frame}

\begin{frame}{Proceso pipeline}
    \centering \textit{Visualización del proceso del pipeline sobre una imagen original}
\end{frame}

\begin{frame}{Segmentación y Morfología}
    \begin{columns}
        \begin{column}{0.5\textwidth}
            \textbf{Segmentación:}
            \begin{itemize}
                \item Umbralización automática mediante \textit{método de Otsu}.
                \item Generación de máscara binaria.
            \end{itemize}
            \textbf{Morfología:}
            \begin{itemize}
                \item Aplicación de \textit{Cierre} para rellenar huecos.
                \item \textit{Erosión/Dilatación} para eliminar ruido.
            \end{itemize}
        \end{column}
        \begin{column}{0.5\textwidth}
            %\includegraphics[width=\textwidth]{segmentacion_placeholder.png}
            \centering \textit{Espacio para visualización de máscara}
        \end{column}
    \end{columns}
\end{frame}



% Slide 5: Análisis Estructural (Bordes y Esquinas)
\begin{frame}{Detección de Bordes y Esquinas}
    \begin{itemize}
        \item \textbf{Bordes:} Implementación del detector de \textit{Canny} para analizar la estructura foliar.
        \item \textbf{Esquinas:} Uso del método de \textit{Shi-Tomasi} para puntos característicos.
        \item \textbf{Discusión:} Análisis de la correlación entre puntos detectados y regiones afectadas por enfermedades.
    \end{itemize}
    \centering \textit{Visualización de bordes sobre imagen original}
\end{frame}

% Slide 6: Descriptores de Características
\begin{frame}{Extracción de Características}
    Para representar las propiedades visuales se utilizaron:
    \begin{block}{Local Binary Patterns (LBP)}
        Captura la información de \textbf{textura} de la hoja.
    \end{block}
    \begin{block}{Histogram of Oriented Gradients (HOG)}
        Captura patrones de \textbf{gradiente y estructura}.
    \end{block}
    \begin{block}{Propuesta Adicional}
        [Insertar aquí su descriptor propuesto o Bag of Words].
    \end{block}
\end{frame}

% Slide 7: Experimentación y Clasificación
\begin{frame}{Estrategia de Clasificación}
    \begin{itemize}
        \item \textbf{Modelos evaluados:} SVM, Random Forest (XGBoost) y k-NN.
        \item \textbf{Validación:} Estrategia k-fold con $k \ge 5$ ante la ausencia de división nativa.
        \item \textbf{Métricas:} Accuracy, F1-score por clase y Matriz de Confusión.
    \end{itemize}
\end{frame}

% Slide 8: Resultados y Conclusiones
\begin{frame}{Resultados Obtenidos}
    \begin{table}[]
        \centering
        \begin{tabular}{lccc}
            \toprule
            Descriptor & Clasificador & Accuracy & F1-Score \\
            \midrule
            LBP        & SVM          & 0.XX     & 0.XX     \\
            HOG        & Random Forest& 0.XX     & 0.XX     \\
            Combinado  & XGBoost      & \textbf{0.XX} & \textbf{0.XX} \\
            \bottomrule
        \end{tabular}
        \caption{Comparativa de desempeño entre descriptores.}
    \end{table}
    \textbf{Conclusión:} La combinación de descriptores [si/no] mejora el desempeño del modelo final.
\end{frame}

% Slide Final
\begin{frame}
    \centering
    \Huge ¿Preguntas?
    \vfill
    \small Visión por Computador - UTEC 2026
\end{frame}

\end{document}